\subsection{Integral finiteness and automorphy for the inertial ring}
\label{sec:rh-dyadic-host-universal}

We now transfer the characteristic-prime curve argument to the
actual ring of the preceding inertial frame.  It is not necessary
to identify this ring with the saturated local-component ring
of the nonresonant argument.

\begin{theorem}[Weight-zero finiteness and the required automorphy]
\label{thm:rh-dyadic-host-universal}
Retain the ordinary RH host and (DH1), with the original admitted
\texttt{HostData}.  There is a mark $g_1$ for the original frame
such that, for every original normalized universal inertial
representation $u$ over its complete local $O$-algebra $R$,
the literal module $\operatorname{wpart}(O,u,0)$ is finite over $O$.
Moreover every type-I point of the original algebra
$\operatorname{Sring}(O,u)$, with coefficients in
$\overline{\mathbf Q}_\ell$, has a torus conjugate which is an
original characteristic-zero point.  Its matrix representation
$\rho$ over $K=\mathbf Q(i)$ remains absolutely irreducible on
every open subgroup and on $G_{K(\mu_{\ell^\infty})}$.
The representation $\rho\epsilon^{-1}$ is attached to a
regular algebraic conjugate-self-dual cuspidal representation
over $K$.

These are exactly the finite-module and automorphy properties
needed in the original \texttt{abstract\_count}, with the
predicate \texttt{Good} chosen to consist of the displayed
automorphy and irreducibility properties.  This assertion does
not assert the additional \texttt{WitA} predicate contained
in the source's stronger \texttt{BPrime\_s}.
\end{theorem}

\begin{proof}
\emph{The common auxiliary frame and the mark.}
The original \texttt{HostFrameEx} construction
(Theorems~\ref{thm:support-auxiliary-tower} and
\ref{thm:support-auxiliary-frame}, with the original seed and
admission propositions) supplies the actual auxiliary field $L$,
its frame, its \texttt{GalLink}, and its second mark.
The hypotheses hold for the present $\psi$:
the auxiliary integer controlling the odd local character orders
can be taken to be one, the $\ell$-power part is trivial, and
the character and decomposition hypotheses were proved in
Theorem~\ref{thm:rh-dyadic-host-frame}.  In particular $L/\mathbf Q$
is abelian and $L=L^+K$, as in (AT1).  The second mark furnished
by this construction has multiplier one and distinct residual
diagonal values, because its image is the corresponding mark
of the auxiliary frame.  Take its image in the inertial
quotient $\Gamma$ for $g_1$.  Use the same complex conjugation,
the same place above $\ell$, and the same marks throughout.

An original inertial point over a local coefficient domain $A$
pulls back along $\pi:G_{\mathbf Q}\to\Gamma$ to an original
\texttt{OrdPt} for this \texttt{GalLink}.  Indeed its multiplier
and residual conditions and both normalizations are the same.
Its \texttt{CondT} is precisely the required ordered local
flag, with a change of basis congruent to one and the original
inertia characters.  Away from $\ell$ it kills $N_v$, and
outside the host set it is unramified.  This is at least the
original \texttt{OrdPt.away} condition, since on inertia away
from $\ell$ the characters $\Psi$ and $\bar\Psi$ have identical
kernels.  The conditions at all conjugates of the chosen
places follow by conjugation.  Adic continuity is preserved
by pullback.  No local saturation or additional equations
are imposed on $A$.

\emph{The actual universal algebra and its weights.}
Write the actual universal ring as
\[
             R=O[[z_1,\ldots,z_m]]/I.                           \tag{DU1}
\]
Such a presentation can be chosen with each $z_i$ an actual
matrix-entry displacement from the original base
representation, hence of weight $w_i\in\{-2,-1,0,1,2\}$.
Here are details of this standard presentation step.
Modulo the square of the maximal ideal and the coefficient
maximal ideal, all entry displacements span the relative
cotangent space: otherwise a nonzero linear functional
annihilating them would give two distinct dual-number maps
with identical universal representations, contradicting
uniqueness in \texttt{IsUniv}.  Choose a finite spanning family
and use complete Nakayama, or successive approximation modulo
powers of the maximal ideal, to obtain the surjection (DU1).

Conjugation by $\operatorname{diag}(t,1,t^{-1})$, for any unit
$t$ in an allowed target ring, preserves the original
normalized inertial functor.  In particular it preserves the
residual diagonal, both normalization equations and the
condition that the local conjugator reduce to one; the
conjugating torus commutes with complex conjugation in the
polarized group.  It preserves \texttt{CondU} as well.
With $t=1+T$ the universal property over $R[[T]]$
therefore realizes the operation $z_i\mapsto(1+T)^{w_i}z_i$.
For a relation in $I$, reduce modulo powers of the maximal
ideal of $R$ and compare its finitely many surviving Laurent
weight components in $1+T$.  Substitution of $1+T$ is
injective on Laurent polynomials over any ring: multiply
by a power of $1+T$ and translate the polynomial variable.
Each weight component of the relation consequently vanishes
in $R$, by separatedness.  Thus $I$ is weight-stable without
a reducedness assumption.

Write $R_j$ for the image of the formal series of weight $j$.
This is exactly the literal $\operatorname{wpart}(O,u,j)$.
One inclusion follows by testing its functional definition
with the universal map to $R[[T]]$ just constructed and
comparing the weight components.  More explicitly, at each
finite maximal-ideal quotient of $R$, any representative
has a finite weight decomposition, and the identity for
$1+T$ forces all components other than weight $j$ to vanish;
passing to the limit gives a weight-$j$ representative.
For the other inclusion, two maps in that functional definition
are local: their targets have the same residue field as $O$,
and a maximal-ideal element cannot map to a nonzero residue
since subtracting an appropriate coefficient unit would send
a unit to a nonunit.  They are consequently continuous.
On the $z_i$ the specified conjugation scales by $t^{w_i}$.
The identity on a homogeneous formal series follows by
continuity, since $t$ and $t^{-1}$ are units.
This checks the original functional definition, not just a
polynomial subalgebra of $R$.

\emph{The characteristic-prime tests.}
Let $f:R\to A$ be any local map to a complete characteristic-
$\ell$ discrete valuation ring with finite residue field.
Specialization gives the actual normalized inertial point:
its original local conjugator specializes, and all prescribed
equations and continuity persist.  Its pullback is the
\texttt{OrdPt} checked above.  The already proved
finite-point theorem and the finite-end transfer
(Theorems~\ref{thm:support-finite-points} and
\ref{thm:support-finite-end}) supply the characteristic
polynomials of the base representation on the fixed
subgroup $G_L$, for every such $A$ and $f$.
For clarity, those results concern the normalized auxiliary
lift of this very point over the same $A$; its inertial
characters on the kernel of the residual prescribed
character are trivial because $A$ has characteristic $\ell$.
Thus its \texttt{FinPt} premise is satisfied.
Lemma~\ref{lem:support-integral-global-triangularity}
now shows that the diagonal entries are the prescribed ones,
and that all the positive or all the negative weight
displacements vanish.  Its proof uses the two original
marks, and applies to this pullback representation without
a change of basis.  We have verified precisely the tests
(IF6) for the actual quotient (DU1).
Proposition~\ref{prop:support-integral-finiteness-criterion}
gives $R_0$ finite over $O$, including the entire original
coefficient torsion and all nilpotents.  By the equality with
the literal weight-zero module, this proves the required
finite-module assertion.

\emph{Lifting points of the original graded algebra.}
Let $\xi:\operatorname{Sring}(O,u)\to
\overline{\mathbf Q}_\ell$ be an original algebra map.
The algebra here consists of finite sums of the $R_j$.
Its restriction to the finite local $O$-algebra $R_0$
has integral values and sends its maximal ideal into
the valuation maximal ideal: every element of the latter
ideal has some power in $\pi_O R_0$, where $\pi_O$ is a
uniformizer.  In particular this assertion applies to the
zero-weight variables and to every nonconstant balanced
monomial.  If $w_i>0>w_j$ and both $\xi(z_i),\xi(z_j)$
are nonzero, the balanced monomial
$z_i^{|w_j|}z_j^{w_i}$ gives
\[
       |w_j|v_\ell(\xi(z_i))+w_i v_\ell(\xi(z_j))>0.           \tag{DU2}
\]
The finitely many inequalities (DU2) guarantee a rational
number $s$ with
$v_\ell(\xi(z_i))+w_i s>0$ for every nonzero $\xi(z_i)$:
each positive-weight lower bound for $s$ is strictly below
each negative-weight upper bound.  If one set is empty,
there is only the other set of bounds.  Choose an algebraic
$t$ with $v_\ell(t)=s$.  Scaling $\xi$ by $t^j$ on weight
$j$ makes all the variable values lie in the maximal ideal
of the integers $O'$ of one finite coefficient extension.
Such an extension exists because only finitely many variable
values and $t$ are involved.

We justify that evaluation at these values extends the scaled
map on every weight space, despite the formal series in (DU1).
For any integer $j$, the set of exponent vectors of
nonnegative integers with total weight $j$ has finitely many
minimal elements in the componentwise ordering, by Dickson's
lemma.  Every such vector lies above one of them, and their
difference has weight zero.  Partition the monomials by a
choice of a minimal vector beneath each.  Every weight-$j$
formal series is therefore a finite sum of those minimal
monomials multiplied by weight-zero formal series.
The map on $R_0$ is continuous: its maximal-ideal and
coefficient-adic topologies are equivalent because it is a
finite local $O$-algebra.  It agrees with evaluation of
weight-zero series by approximating them by polynomials in
the finitely generated monoid of balanced monomials, whose
nonconstant generators lie in its maximal ideal.
The finite weight-space decomposition just proved makes
the same agreement hold at every weight.  For any relation
in $I$, all its homogeneous parts belong to $I$.  Evaluation
of each part is zero.  Regrouping the convergent formal series
by weight is valid because all variable values are strictly
small.  Hence every relation evaluates to zero and we
obtain a local homomorphism $R\to O'$.

Specialize $u$ along this map.  By the agreement on all its
entries and the defining identity \texttt{IsPtHom}, the
resulting representation is precisely the torus conjugate
of the original point represented by $\xi$.
It satisfies the original normalized local and global
conditions and is adically continuous.  This gives its
literal \texttt{IsCharZeroPt} witness, with the same
representation up to the permitted torus conjugation,
not a new point on a different ring.

\emph{Automorphy of a type-I point.}
Type I is absolute irreducibility on the multiplier-one
subgroup.  The original ordinary flag and the separating
decomposition-group element (DH3) satisfy the hypotheses
of Lemma~\ref{lem:support-type-one-strong-irreducibility}.
Thus the matrix representation $\rho$ has all the
irreducibility properties in the statement, including
irreducibility after restriction to $G_L$.
The actual normalized-lift transfer (FE1) over the coefficient
DVR and the original \texttt{GoodPt} theorem give the
arithmetic point of the actual auxiliary Hecke quotient,
with the determinant of $\rho|_{G_L}\epsilon^{-1}$ and
its exact ordered inertia characters.
The argument (TR4) in
Theorem~\ref{thm:support-type-one-rigidity}, using the
proved arithmetic-point classicality for that quotient,
therefore gives a regular algebraic conjugate-self-dual
cuspidal representation of $\mathrm{GL}_3(\mathbf A_L)$
attached to $\rho|_{G_L}\epsilon^{-1}$.
Every step uses the original normalized lift; it does not
require the inertial ring to be the saturated ring.

Apply soluble CM descent
(Thorne, \emph{Automorphy lifting for residually reducible
$l$-adic Galois representations}, Lemma 2.7(2))
to $\rho\epsilon^{-1}$ over $K$.
Indeed $L/K$ is abelian, the representation on $G_L$
is irreducible, and its conjugate-dual multiplier is the
restriction of $\epsilon^{-2}$ on $\mathbf Q$, with
constant value on complex conjugation.  Descent produces
a regular algebraic cuspidal representation over $K$
whose Galois representation is the stated one.
Its conjugate self-duality in the automorphic normalization
follows from the original polarized identity with multiplier
$\epsilon^{-2}$ and strong multiplicity one at the
unramified places.  The Tate twist and its sign have
not changed.  This supplies exactly the automorphy and
irreducibility needed for the preceding rigidity and
Frobenius-spectrum arguments.
\end{proof}

\paragraph{Why the stronger automorphy witness is unnecessary.}
The original \texttt{abstract\_count} is quantified over an
arbitrary predicate \texttt{Good}.  Its use is only that every
type-I graded point has this predicate, and that its
characteristic-zero type-I points with the predicate are rigid
and have nonzero trace difference.  Take \texttt{Good} to mean
the irreducibility and automorphy just proved.  The proofs of
Theorem~\ref{thm:rh-dyadic-host-traces} use only these properties
of \texttt{BPrime\_s}: the local crystallinity proof and
the spectrum theorem replace the separate local witness.
In particular no assertion about the nonresonance of
all powers of the diagonal character ratios in
\texttt{WitA} is needed or inferred.  This is a logically
sufficient weaker interface for the original integral
counting theorem, not a change to its arithmetic points
or local conditions.

The remaining step at three in this route is the
integral graded-algebra estimate for the actual inertial
universal ring, including the adapted presentation and its
exact coefficient-torsion contributions.  Once that estimate
is proved, the arithmetic inputs of the count, the finite
module, the point automorphy, rigidity, strict length
comparison and uniform trace bound are all available.
The full W04 node remains open.
